\section{Introduction}

\subsection{Importance and Relevance to Computer Networks and Society}
Computer networking constitutes the foundational digital arterial system of modern human civilization \cite{kurose2021computer}. From cloud-native enterprise operations and financial transaction clearinghouses to critical healthcare data systems, modern society is utterly reliant upon high-speed, dependable, and secure packet transmission \cite{tanenbaum2021computer}. The design, optimization, and security fortification of transport layers, switching fabrics, and routing backbones determine the efficiency and resilience of global communications.

In this case study, we investigate \textbf{Network Security Analysis: Firewall and Intrusion Detection Systems}, contextualized within enterprise campus network comprising 1,800 host endpoints, DMZ web services, and internal database servers subject to continuous external scanning and volumetric attacks. As network throughput requirements escalate from gigabit to terabit scales, network engineers must resolve complex tradeoffs between latency, security inspection depth, protocol overhead, and computational resource consumption. A single misconfiguration in routing protocols or a vulnerability in perimeter firewall infrastructure can trigger catastrophic outages or expose sensitive institutional assets to adversarial exploitation.

\subsection{Objectives of the Case Study Analysis}
The fundamental engineering and investigative objectives of this study are:
\begin{itemize}
  \item To establish the rigorous theoretical and mathematical foundations governing protocol state machines, packet encapsulation, and network algorithms for Network Security Analysis: Firewall and Intrusion Detection Systems.
  \item To architect a robust, production-grade network topology addressing operational bottlenecks: High false positive alert volume in signature engines, state table exhaustion from SYN flood denial-of-service attacks, and inspection latency during 10 Gbps line-rate forwarding..
  \item To conduct empirical performance benchmarking using industry-standard tools (Wireshark, iperf3, Snort, GNS3), measuring packet delivery, latency, and mitigation efficacy.
  \item To formulate systematic diagnostic workflows for rapid packet-level troubleshooting and protocol optimization.
  \item To map the technical outcomes directly to ABET Program Outcomes and the United Nations Sustainable Development Goals.
\end{itemize}
